% Options for packages loaded elsewhere
\PassOptionsToPackage{unicode}{hyperref}
\PassOptionsToPackage{hyphens}{url}
\PassOptionsToPackage{dvipsnames,svgnames,x11names}{xcolor}
\documentclass[
  10pt,
  fontset=fandol]{ctexart}
\usepackage{xcolor}
\usepackage[margin=16mm]{geometry}
\usepackage{amsmath,amssymb}
\setcounter{secnumdepth}{-\maxdimen} % remove section numbering
\usepackage{iftex}
\ifPDFTeX
  \usepackage[T1]{fontenc}
  \usepackage[utf8]{inputenc}
  \usepackage{textcomp} % provide euro and other symbols
\else % if luatex or xetex
  \usepackage{unicode-math} % this also loads fontspec
  \defaultfontfeatures{Scale=MatchLowercase}
  \defaultfontfeatures[\rmfamily]{Ligatures=TeX,Scale=1}
\fi
\usepackage{lmodern}
\ifPDFTeX\else
  % xetex/luatex font selection
\fi
% Use upquote if available, for straight quotes in verbatim environments
\IfFileExists{upquote.sty}{\usepackage{upquote}}{}
\IfFileExists{microtype.sty}{% use microtype if available
  \usepackage[]{microtype}
  \UseMicrotypeSet[protrusion]{basicmath} % disable protrusion for tt fonts
}{}
\makeatletter
\@ifundefined{KOMAClassName}{% if non-KOMA class
  \IfFileExists{parskip.sty}{%
    \usepackage{parskip}
  }{% else
    \setlength{\parindent}{0pt}
    \setlength{\parskip}{6pt plus 2pt minus 1pt}}
}{% if KOMA class
  \KOMAoptions{parskip=half}}
\makeatother
\setlength{\emergencystretch}{3em} % prevent overfull lines
\providecommand{\tightlist}{%
  \setlength{\itemsep}{0pt}\setlength{\parskip}{0pt}}
\usepackage{amsmath,amssymb,mathtools,needspace}
\xeCJKDeclareCharClass{CJK}{"2032,"0400->"04FF,"2160->"216B,"2460->"2473,"25A0,"25C7,"25CB,"25CF}
\IfFontExistsTF{Arial Unicode MS}{\xeCJKsetup{AutoFallBack=true}\setCJKfallbackfamilyfont{\CJKrmdefault}{Arial Unicode MS}}{}
\setcounter{secnumdepth}{0}
\setlength{\emergencystretch}{2em}
\allowdisplaybreaks
\usepackage{bookmark}
\IfFileExists{xurl.sty}{\usepackage{xurl}}{} % add URL line breaks if available
\urlstyle{same}
\hypersetup{
  pdftitle={正交化手续},
  colorlinks=true,
  linkcolor={Maroon},
  filecolor={Maroon},
  citecolor={Blue},
  urlcolor={Blue},
  pdfcreator={LaTeX via pandoc}}

\title{正交化手续}
\author{}
\date{}

\begin{document}
\maketitle

\subsubsection{3.4.1 正交化手续}\label{ux6b63ux4ea4ux5316ux624bux7eed}

定义3.9 设 \(g_n(x)\) 是首项系数 \(a_n \neq 0\) 的 \(n\) 次多项式，如果多项式序列 \(g_0(x)\), \(g_1(x)\), \(\cdots\) 满足 \[
(g_j, g_k) = \int_a^b \rho(x) g_j(x) g_k(x) \mathrm{d}x = \begin{cases} 0, & j \neq k, \\ A_k > 0, & j = k \end{cases} \quad (j, k = 0, 1, \cdots),
\] 则称多项式序列 \(g_0(x), g_1(x), \cdots\) 在 \([a, b]\) 上带权 \(\rho(x)\) 正交，并称 \(g_n(x)\) 是 \([a, b]\) 上带权 \(\rho(x)\) 的 \(n\) 次正交多项式.

一般来说，当权函数 \(\rho(x)\) 及区间 \([a, b]\) 给定以后，可以由线性无关的一组基 \(\{1,\allowbreak  x,\allowbreak  x^2,\allowbreak  \cdots,\allowbreak  x^n,\allowbreak  \cdots\}\) 并利用正交化方法构造出正交多项式 \[
g_0(x) = 1, \quad g_n(x) = x^n - \sum_{k=0}^{n-1} \frac{(x^n, g_k)}{(g_k, g_k)} \cdot g_k(x), \quad k = 1, 2, \cdots.
\] 这样构造的正交多项式有以下性质： (1) \(g_n(x)\) 是最高项系数为 1 的 \(n\) 次多项式. (2) 任一 \(n\) 次多项式 \(P_n(x) \in H_n\) 均可表示为 \(g_0(x), g_1(x), \cdots, g_n(x)\) 的线性组合. (3) 当 \(n \neq m\) 时，\((g_n, g_m) = 0\) 且 \(g_n(x)\) 与任一次数小于 \(n\) 的多项式正交. (4) 有递推关系 \[
g_{n+1}(x) = (x - \alpha_n)g_n(x) - \beta_n g_{n-1}(x), \quad n = 0, 1, \cdots,
\] 其中 \[
g_0(x) = 1, \quad g_{n-1}(x) = 0,
\] \[
\alpha_n = \frac{(xg_n, g_n)}{(g_n, g_n)}, \quad n = 0, 1, \cdots,
\] \[
\beta_n = \frac{(g_n, g_n)}{(g_{n-1}, g_{n-1})}, \quad n = 1, 2, \cdots.
\] 这里 \((xg_n, g_n) = \int_a^b xg_n^2(x) \rho(x) \mathrm{d}x.\) (5) 设 \(g_0(x), g_1(x), \cdots\) 是在 \([a, b]\) 上带权 \(\rho(x)\) 的正交多项式序列，则 \(g_n(x)\) (\(n \geq 1\)) 的 \(n\) 个根都是单重实根，且都在区间 \((a, b)\) 内. 以上性质的证明可参看相关文献. 下面主要给出几类常见的而又十分重要的正交多项式.

\begin{center}\rule{0.5\linewidth}{0.5pt}\end{center}

\subsection{相关原页校注（agent 补充）}\label{ux76f8ux5173ux539fux9875ux6821ux6ce8agent-ux8865ux5145}

以下校注来自本节覆盖的原页，可能也涉及同页相邻小节；原文照录与校正说明分开保留。

\subsubsection{原书 PDF 页 70 的校注}\label{ux539fux4e66-pdf-ux9875-70-ux7684ux6821ux6ce8}

\href{../pages/pdf-070.md}{完整原页转录} · \href{../../source-images/pdf-070.jpeg}{原页图像}

校注记录：ch03a-E08, ch03a-E09。

\begin{quote}
校注（agent 补充，原书索引疑误）：本页正交化构造式末尾原印``\(k=1,2,\cdots\)''，而 \(k\) 已是求和哑指标，外部序号应为 \(n=1,2,\cdots\)；转录保留原式。
\end{quote}

\begin{quote}
校注（agent 补充，原书初值疑误）：递推初值原印 \(g_{n-1}(x)=0\)，按原文保留。若对一般 \(n\) 成立，会与 \(g_0(x)=1\) 冲突；递推通常应以 \(g_{-1}(x)=0\) 开始。
\end{quote}

\subsection{本节覆盖原页的校注记录}\label{ux672cux8282ux8986ux76d6ux539fux9875ux7684ux6821ux6ce8ux8bb0ux5f55}

下列记录按原页关联，含同页相邻小节的校注；计算时同时读取上文校注和完整原页。

\begin{itemize}
\tightlist
\item
  \texttt{ch03a-E08}：原书 PDF 页 70
\item
  \texttt{ch03a-E09}：原书 PDF 页 70
\end{itemize}

\end{document}
